\scanpage{118}
\deleted{Page 33. Ligne -9 à -1, remplacer le texte par le suivant:

« \uline{Remarques} 7.10. a) Lorsque $G$ est à fibres connexes, alors on peut dans c) abandonner l'hypothèse que les \deleted{$\unsure{p}$} $\inserted{{}_{\unsure{p^\nu|p^n}}}H'$ soient séparés. En effet, avec les réductions faites dans la démonstration, on peut supposer $u$ bijectif, donc $H$ à fibres connexes, donc $H$ séparé en vertu du théorème de séparation de Raynaud (Exp $\mathrm{VI}_{\mathrm B}$\quad).

b) D'autre part, des exemples de Raynaud montrent qu'on ne peut en général dans c) abandonner ni l'hypothèse que les $\inserted{{}_{\unsure{p^\nu|p^n}}}G'$ soient finis sur $S'$, ni que les $\inserted{{}_{\unsure{p^\nu|p^n}}}H'$ soient séparés sur $S'$.

c) \inserted{On peut montrer que} \unsure{La|la} conclusion de a) est encore valable si, au lieu de supposer $G$ à fibres connexes, on le suppose de type multiplicatif; \deleted{\unsure{alors}} tout monomorphisme $u:G\to H$ de $S$-préschémas en groupes\deleted{\unsure{à}} de présentation finie, avec $G$ de type multiplicatif, \deleted{\unsure{u}} est une immersion. »

Ligne 14, lire $S'$ au lieu de $S$.}

Page 35, ligne 9, 10, supprimer ``lisse sur $k$''.
\mspage{117}
